Key human studies on “leaky gut” (increased intestinal permeability)
“Leaky gut” refers to impaired intestinal barrier function, typically measured by dual-sugar tests (e.g., lactulose/mannitol or lactulose/rhamnose ratio), urinary recovery of probes (sucralose, etc.), or circulating biomarkers (zonulin, LPS/endotoxin, LBP, I-FABP/FABP2). Below is a selection of important human studies and reviews, grouped by category. These focus on measurement, associations with disease, and interventions.
Foundational / Measurement Studies
O’Dwyer et al. (1988): Single IV dose of E. coli endotoxin (4 ng/kg) in 12 healthy volunteers approximately doubled lactulose absorption/excretion versus saline, demonstrating that circulating endotoxin rapidly increases intestinal permeability.
Camilleri (2019) review (Gut): Comprehensive overview of mechanisms, in vivo (sugar probes) and in vitro measurement methods, effects of stressors (NSAIDs, exercise, pregnancy), and dietary modulators in humans. Emphasizes that while barrier dysfunction occurs in many conditions, normalizing it alone does not cure diseases.
Khoshbin et al. / Mayo Clinic group (2021): Validated multi-sugar test (¹³C-mannitol, rhamnose, lactulose, sucralose) in 60 healthy adults for small-intestinal vs colonic permeability. Provided normative data; ¹³C-mannitol preferred for small-bowel assessment; fiber intake had minor effects. Pilot data in IBS-D.
Disease Associations
Forsyth et al. (2011): Early Parkinson’s disease patients showed significantly higher 24-hour urinary sucralose (total gut permeability) versus age-matched controls; correlated with mucosal E. coli, nitrotyrosine, α-synuclein staining, and serum LBP.
Stevens et al. (2018): Elevated plasma zonulin and FABP2 correlated with plasma LPS and altered gut microbiome in patients with anxiety or depression.
Multiple cardiovascular/endotoxemia cohorts (e.g., Hisayama Study on serum LBP and incident CVD; zonulin/LBP/I-FABP predicting cardiovascular mortality over 15 years; higher LPS/zonulin post-MI linked to outcomes and smoking). Systematic review/meta-analysis (2024) confirmed higher intestinal permeability markers (LPS, zonulin, I-FABP, etc.) in cardiovascular disease patients versus controls.
Centenarian vs early MI comparison: Disease-free centenarians had lower serum zonulin and LPS than young patients with precocious myocardial infarction (and lower LPS than young healthy controls).
Intervention / Treatment Trials
Lubiprostone (2017 pilot RCT): 24 µg/day for 28 days in healthy men reduced NSAID (diclofenac)-induced increase in lactulose-mannitol ratio versus untreated controls—the first reported improvement of “leaky gut” with an available drug in humans.
Multistrain probiotics:
Pilot study in IBS-D patients with confirmed increased permeability: 30 days of multistrain probiotic improved permeability in ~81% (normalized in 37%), with better stool consistency and quality of life.
RCT in elderly Thai subjects: Multi-strain (L. paracasei, B. breve, B. longum) for 12 weeks reduced lactulose-mannitol ratio (improved barrier up to ~48%), lowered LPS/hs-CRP, and improved metabolic markers.
2025 systematic review & meta-analysis of RCTs: Pro-/synbiotics significantly reduced LPS (SMD −0.54) and zonulin (SMD −0.49); prebiotics reduced LPS (SMD −0.88). Supports barrier reinforcement, though heterogeneity is high.
Other interventions: Glutamine, zinc carnosine, vitamin D (in Crohn’s), and certain fibers/oats have shown benefits on permeability markers in smaller human trials; results for probiotics overall are mixed in some meta-analyses.
Additional Notes
Common stressors that increase permeability in humans: NSAIDs, intense exercise, critical illness, western diet components (emulsifiers, high fructose), and metabolic conditions (obesity, metabolic syndrome, type 2 diabetes).
Biomarkers such as zonulin, LPS, LBP, and I-FABP are frequently used alongside (or instead of) classic sugar-probe tests; they correlate with each other and with systemic inflammation/cardiovascular risk in multiple cohorts.
“Leaky gut syndrome” as a standalone diagnosis remains controversial in mainstream medicine; evidence is strongest for secondary barrier dysfunction in specific diseases (IBD, IBS, metabolic disorders, neurodegenerative conditions, critical illness) rather than as a primary cause of unrelated systemic illnesses.
These studies form the core human evidence base. Larger, longer-term RCTs with standardized permeability assays are still needed for many proposed interventions.
